\section*{Введение}
\addcontentsline{toc}{section}{Введение}

Объём научной литературы растёт экспоненциально: по разным оценкам число публикаций
удваивается примерно каждые 10–15 лет~\cite{bornmann2015growth}. На таком фоне исследователю
всё труднее удерживать в голове карту собственной области — отслеживать смежные направления,
видеть, где его работа граничит с чужой, и находить людей, с которыми имеет смысл начать
совместный проект. Классические библиометрические инструменты опираются прежде всего на графы
цитирований и соавторства~\cite{newman2001coauthorship}. Такие графы хорошо описывают уже
состоявшиеся связи, но почти не подсказывают новые: задача предсказания ещё не возникших рёбер
(link prediction) в социальных и научных сетях известна своей сложностью~\cite{libennowell2007link},
а главное — формальная связь <<кто кого процитировал>> слабо отражает смысловую близость работ.

Ситуация изменилась с приходом трансформерных архитектур. Сначала
BERT~\cite{devlin2019bert} показал, что предобученный энкодер хорошо улавливает контекст слова,
а затем энкодерные модели эмбеддингов вроде Sentence-BERT~\cite{reimers2019sbert} научились
кодировать смысл целого текста в один вектор так, что семантически близкие документы
оказываются рядом в пространстве. Это позволяет перейти от формальных связей к семантическому
анализу — представить статью вектором, сгруппировать корпус по смыслу, а не по совпадению
ключевых слов, и оценить близость двух исследователей по содержанию их работ, даже если в
списке соавторов они никогда не пересекались. Параллельно большие языковые
модели (LLM) дали возможность не только ранжировать кандидатов, но и объяснять рекомендацию на
естественном языке, а заодно выступать оценщиком качества (LLM-as-a-judge) там, где ручная
разметка слишком дорога~\cite{zheng2023judging}. Наконец, открытые базы вроде
OpenAlex~\cite{priem2022openalex} сделали корпус из сотен тысяч статей с метаданными доступным
без коммерческих ограничений.

\textbf{Тема работы} — система кластеризации научных статей.

\textbf{Объект исследования} — научное информационное пространство: корпус публикаций и
исследователи, стоящие за ними. \textbf{Предмет исследования} — методы векторного представления
статей, тематической кластеризации публикаций и поиска семантически совместимых соавторов с
генерацией объяснений.

Отдельно поясню постановку задачи рекомендаций — здесь легко допустить неточность.
Рекомендатель состоит из двух частей. Первая, обучаемая, — этап извлечения (retrieval):
нейросетевая модель по истории публикаций автора строит вектор его научных интересов, и мы
отбираем авторов с наиболее близкими векторами. То есть retrieval находит исследователей с
близкой тематической траекторией, а не «дополняющих» друг друга — это важно не путать. Вторая
часть — переранжирование на предобученной языковой модели: она оценивает релевантность
отобранных кандидатов и уже способна выделять взаимодополняющую экспертизу, формируя
объяснение. Обучается в системе только первая часть; языковая модель используется без
дообучения. Результат мы сознательно не называем «предсказанием будущего соавтора»: фактические
соавторства — разреженный и неполный сигнал, они покрывают лишь часть всех релевантных связей,
поэтому корректнее говорить об извлечении и ранжировании совместимых кандидатов (количественно
этот разрыв показан в главе~3).

\textbf{Цель работы} — разработать систему семантического анализа научного пространства, которая
строит векторные представления статей, выполняет интерпретируемую тематическую кластеризацию,
визуализирует пространство в виде интерактивной карты и рекомендует совместимых соавторов с
понятными объяснениями.

Для достижения цели поставлены следующие \textbf{задачи}:

\begin{enumerate}
  \item собрать корпус научных публикаций из OpenAlex и привести тексты к единому виду;
  \item построить многоязычные эмбеддинги статей;
  \item реализовать пайплайн тематической кластеризации в стиле BERTopic;
  \item сравнить HDBSCAN с наивными бейзлайнами (K-Means и случайное разбиение на UMAP-10)
        по набору независимых метрик и обоснованно выбрать лучший вариант;
  \item построить семантический слой метакластеров — укрупнённое объединение мелких тем в крупные
        группы;
  \item обучить нейросетевую модель представления автора и реализовать retrieval кандидатов;
  \item реализовать LLM-переранжирование кандидатов и генерацию объяснений;
  \item разработать веб-приложение с интерактивной картой и интерфейсом рекомендаций;
  \item провести эксперименты: оценить качество кластеризации, retrieval и рекомендаций.
\end{enumerate}

\textbf{Методы исследования.} Векторизация статей — multilingual-e5-base~\cite{wang2024multilinguale5};
понижение размерности — UMAP~\cite{mcinnes2018umap}; кластеризация — HDBSCAN~\cite{campello2013hdbscan,
mcinnes2017hdbscanjoss}; интерпретируемое представление тем — c-TF-IDF
и Maximal Marginal Relevance~\cite{carbonell1998mmr} в духе BERTopic~\cite{grootendorst2022bertopic};
представление автора — Transformer-энкодер~\cite{vaswani2017attention} и GraphSAGE
поверх графа соавторства~\cite{hamilton2017graphsage} с обучением по coauthor InfoNCE;
оценка качества — proxy-метрики по будущим соавторам, в том числе nDCG~\cite{jarvelin2002ndcg},
и LLM-as-a-judge~\cite{zheng2023judging} для релевантности финальной выдачи.

\textbf{Основной результат.} Реализована работающая система <<AI Coauthor>>. На корпусе из
509\,447 статей по подполю Artificial Intelligence и 35\,954 квалифицированных авторов построены
979 тем и 76 метакластеров. HDBSCAN выбран по совокупности метрик качества тем с явным
компромиссом по доле шума. Рекомендательный пайплайн
сравнен в трёх режимах — mean-бейзлайн, Transformer и GraphSAGE; LLM Has3@10 (score $\geq 2$)
равен 0,875, а LLM Selected3Mean@10 достигает 2,49 из 3. Графовый метод даёт
сопоставимое качество с Transformer, но не показывает значимого выигрыша от кластерных узлов. Все
подсистемы объединены в веб-приложение
с интерактивной картой научного пространства.

\textbf{Структура работы.} Работа состоит из введения, трёх глав, заключения и
библиографического списка. В главе~1 дан обзор предметной области: научные графы, текстовые
эмбеддинги, тематическое моделирование, рекомендация соавторов и подход LLM-as-a-judge. В
главе~2 описаны архитектура и реализация системы AI~Coauthor — от сбора данных и построения
эмбеддингов до кластеризации, метакластеров, рекомендательного пайплайна и веб-приложения. В
главе~3 приведены эксперименты: описание датасета, сравнение кластеризации с бейзлайнами,
retrieval-метрики и LLM-оценка качества рекомендаций. В заключении подведены итоги и обозначены
ограничения и направления развития.
